%! Introduction to Eigenvalues — Unit 6, Lesson 6.1 — Student Packet Cover Sheet
\documentclass[10pt]{article}
\usepackage{linalg-article}
\usepackage{linalg-boxes}
\usepackage{ltablex}
\keepXColumns

\begin{document}

% Full-bleed burgundy banner
\begin{tikzpicture}[remember picture, overlay]
  \fill[burgundy] (current page.north west)
    rectangle ([yshift=-0.9in]current page.north east);
\end{tikzpicture}
\vspace{-0.8in}

\begin{center}
  {\color{white}\textbf{\LARGE Linear Algebra}}\\[4pt]
  {\color{blushmid}\large\textbf{Unit 6 \quad Eigenvalues and Eigenvectors}}\\[6pt]
  {\color{white}\textbf{\large Lesson 6.1 \quad Introduction to Eigenvalues}}
\end{center}

\vspace{0.22in}
\namedateperiod
\vspace{0.18in}

\begin{learningtargetbox}
\small
\begin{enumerate}[leftmargin=1.5em, itemsep=5pt]
  \item \ldots{} turn $A\mathbf{x}=\lambda\mathbf{x}$ into $(A-\lambda I)\mathbf{x}=\mathbf{0}$ and explain why the eigenvalues satisfy $\det(A-\lambda I)=0$.
  \item \ldots{} \textbf{find} a matrix's eigenvalues by solving the characteristic equation $\det(A-\lambda I)=0$.
  \item \ldots{} \textbf{find} an eigenvector for each $\lambda$ by solving $(A-\lambda I)\mathbf{x}=\mathbf{0}$, and check with the trace and determinant.
\end{enumerate}
\end{learningtargetbox}

\vspace{0.14in}

\begin{tocbox}
\small
\renewcommand{\arraystretch}{1.5}
\begin{tabularx}{\linewidth}{@{} c l X r @{}}
\rowcolor{blush}
\textbf{\#} & \textbf{Component} & \textbf{Description} & \textbf{Score} \\
\midrule
1 & Warm-Up        & Spiral review: $A-\lambda I$, the $2\times2$ determinant, solving a quadratic & \blank{1.2cm} \\
2 & Guided Notes   & The characteristic equation $\det(A-\lambda I)=0$; finding $\lambda$'s and eigenvectors & \blank{1.2cm} \\
3 & Group Activity & Running the full method; a $\lambda=0$ matrix; a rotation with no real $\lambda$ & \blank{1.2cm} \\
4 & Exit Ticket    & Quick check: find eigenvalues, an eigenvector, and justify the method            & \blank{1.2cm} \\
5 & Homework       & Full method, diagonal/triangular shortcut, a negative $\lambda$, trace/det check   & \blank{1.2cm} \\
\midrule
  & \multicolumn{2}{r}{\textbf{Total}} & \blank{1.2cm} \\
\end{tabularx}
\end{tocbox}

\vspace{0.10in}

\begin{remindbox}
\small Yesterday you \emph{tested} vectors you were handed. Today you \emph{find} the special
directions from scratch. The whole method is two steps: solve the \textbf{characteristic equation}
$\det(A-\lambda I)=0$ for the \textbf{eigenvalues} $\lambda$, then solve $(A-\lambda I)\mathbf{x}=
\mathbf{0}$ for an \textbf{eigenvector} of each. A nonzero $\mathbf{x}$ can be crushed to $\mathbf{0}$
only when $A-\lambda I$ is singular --- that is why the determinant must be zero (Unit~5). Two fast
checks: the eigenvalues \emph{sum} to the trace and \emph{multiply} to the determinant.
\end{remindbox}

\end{document}
